\documentclass[11pt,a4paper]{ctexart}
\usepackage{amsmath,amssymb,mathtools,bm}
\usepackage{geometry}
\usepackage{enumitem}
\usepackage{hyperref}
\usepackage{xcolor}
\usepackage{fancyhdr}
\usepackage{titlesec}
\usepackage{booktabs,longtable,array}

\geometry{left=2.05cm,right=2.05cm,top=2.0cm,bottom=2.2cm}
\setmonofont{Consolas}
\hypersetup{
  colorlinks=true,
  linkcolor=blue!45!black,
  urlcolor=blue!45!black,
  citecolor=blue!45!black,
  pdftitle={数值代数期末备考详细自学讲义},
  pdfauthor={Codex}
}
\setlist{nosep,leftmargin=2em}
\allowdisplaybreaks
\sloppy
\pagestyle{fancy}
\setlength{\headheight}{14pt}
\fancyhf{}
\fancyhead[L]{数值代数期末备考}
\fancyhead[R]{数值代数期末备考详细自学讲义}
\fancyfoot[C]{\thepage}
\renewcommand{\headrulewidth}{0.4pt}
\newcommand{\code}[1]{\begingroup\ttfamily\small #1\endgroup}
\titlespacing*{\section}{0pt}{1.4em}{0.5em}
\titlespacing*{\subsection}{0pt}{1.0em}{0.35em}
\titlespacing*{\subsubsection}{0pt}{0.8em}{0.25em}

\title{\bfseries 数值代数期末备考详细自学讲义}
\author{基于历年卷、非上机作业与课程讲义整理}
\date{2026年6月26日}

\begin{document}
\maketitle
\tableofcontents
\newpage

\phantomsection

\section*{使用说明}

\addcontentsline{toc}{section}{使用说明}

这份讲义按“从忘得差不多到能上考场”的标准写。读的时候不要只背结论，建议每一章都按下面的顺序走：

\begin{enumerate}

\item 先看“本章解决什么问题”，知道这章在整门课中的位置；

\item 再看定义和核心定理，理解每个符号是什么意思；

\item 跟着算法步骤做一遍小矩阵手算；

\item 最后背“考场模板”和“易错点”。

\end{enumerate}

如果时间很紧，优先顺序是：

\begin{enumerate}

\item LU/PLU、Cholesky/LDL\textasciicircum{}T、最小二乘、Householder/Givens；

\item 范数、条件数、投影矩阵、证明题模板；

\item Jacobi/G-S/SOR、CG；

\item Hessenberg、QR、Schur、SVD、幂法、二分法。

\end{enumerate}

整门课可以看成四句话：

\begin{itemize}

\item 解线性方程组：把困难矩阵变成三角矩阵；

\item 解最小二乘：把残差问题变成正交投影或 QR；

\item 解迭代问题：把误差写成 $e^{(k+1)}=Be^{(k)}$，然后看 $\rho(B)<1$；

\item 解特征值问题：用正交相似变换保持特征值，同时逐步逼近三角形结构。

\end{itemize}

\phantomsection

\section*{零、必备线性代数与记号}

\addcontentsline{toc}{section}{零、必备线性代数与记号}

\phantomsection

\subsection*{1. 向量、矩阵和基本空间}

\addcontentsline{toc}{subsection}{1. 向量、矩阵和基本空间}

数值代数关心的是“如何稳定、有效地计算”。但所有算法都建立在线性代数语言上。

线性方程组写作：

\[
Ax=b,\quad A\in\mathbb R^{m\times n},\quad x\in\mathbb R^n,\quad b\in\mathbb R^m.
\]

当 $m=n$ 且 \code{\detokenize{A}} 非奇异时，方程组有唯一解：

\[
x=A^{-1}b.
\]

考试中很少真的让你求 $A^{-1}$ 再乘 \code{\detokenize{b}}，因为这样数值上不经济。真正常用的是：

\begin{enumerate}

\item 把 \code{\detokenize{A}} 分解成容易解的矩阵；

\item 分两次或多次解三角方程；

\item 必要时用正交变换控制误差。

\end{enumerate}

常用空间：

\begin{itemize}

\item \code{\detokenize{R(A)}}：矩阵 \code{\detokenize{A}} 的列空间，也就是所有 \code{\detokenize{Ax}} 组成的空间；

\item \code{\detokenize{N(A)}}：矩阵 \code{\detokenize{A}} 的零空间，也就是所有满足 $Ax=0$ 的向量；

\item \code{\detokenize{rank(A)}}：列空间维数，也等于行空间维数；

\item \code{\detokenize{A}} 满列秩：$rank(A)=n$，列向量线性无关；

\item \code{\detokenize{A}} 满秩方阵：$rank(A)=n$，等价于 \code{\detokenize{A}} 可逆。

\end{itemize}

一个非常常用的结论是：

\[
N(A^TA)=N(A).
\]

证明很短，但考试很爱用。若 $Ax=0$，显然 $A^TAx=0$。反过来，若 $A^TAx=0$，左乘 $x^T$ 得

\[
0=x^TA^TAx=\|Ax\|_2^2,
\]

因此 $Ax=0$。这个结论的用处是：如果 \code{\detokenize{A}} 满列秩，那么 $A^TA$ 非奇异，最小二乘正则方程有唯一解。

\phantomsection

\subsection*{2. 对称正定矩阵}

\addcontentsline{toc}{subsection}{2. 对称正定矩阵}

对称正定矩阵是本课程的一个中心对象。定义是：

\[
A=A^T,\quad x^TAx>0\quad (x\ne 0).
\]

它经常简写为 SPD。

等价判别方法：

\begin{enumerate}

\item 所有特征值都大于 0；

\item 所有顺序主子式都大于 0；

\item 能做 Cholesky 分解且所有对角元为正；

\item 二次型 $x^TAx$ 可以配方成正平方和。

\end{enumerate}

顺序主子式判别的意思是：对称矩阵 \code{\detokenize{A}} 正定当且仅当

\[
\det A_1>0,\quad \det A_2>0,\quad \cdots,\quad \det A_n>0,
\]

其中 $A_k$ 是 \code{\detokenize{A}} 左上角 \code{\detokenize{k}} 阶主子矩阵。

SPD 的三个重要后果：

\begin{itemize}

\item 可以做 Cholesky；

\item 共轭梯度法适用；

\item 对任意非零 \code{\detokenize{x}}，$x^TAx$ 可以作为一种“能量”或“A-范数”的平方。

\end{itemize}

\phantomsection

\subsection*{3. 正交矩阵}

\addcontentsline{toc}{subsection}{3. 正交矩阵}

正交矩阵满足：

\[
Q^TQ=QQ^T=I.
\]

如果是复矩阵，对应的是酉矩阵：

\[
Q^*Q=QQ^*=I.
\]

正交矩阵最重要的性质是保持 2-范数：

\[
\|Qx\|_2^2=(Qx)^T(Qx)=x^TQ^TQx=x^Tx=\|x\|_2^2.
\]

所以 Householder、Givens、QR、Hessenberg 化、QR 迭代都喜欢用正交变换。考场上看到“为什么稳定”“为什么最小二乘可以这样化简”“为什么特征值不变”，多半都要写：

\[
Q^TQ=I.
\]

\phantomsection

\subsection*{4. 三角矩阵为什么重要}

\addcontentsline{toc}{subsection}{4. 三角矩阵为什么重要}

如果 \code{\detokenize{L}} 是下三角，解 $Ly=b$ 可以前代；如果 \code{\detokenize{U}} 是上三角，解 $Ux=y$ 可以回代。

以上三角为例：

\[
\begin{pmatrix}
u_{11}&u_{12}&\cdots&u_{1n}\\
0&u_{22}&\cdots&u_{2n}\\
\vdots&\vdots&\ddots&\vdots\\
0&0&\cdots&u_{nn}
\end{pmatrix}
\begin{pmatrix}x_1\\x_2\\\vdots\\x_n\end{pmatrix}
=
\begin{pmatrix}y_1\\y_2\\\vdots\\y_n\end{pmatrix}.
\]

回代从最后一行开始：

\[
x_n=\frac{y_n}{u_{nn}},
\]

然后

\[
x_i=\frac{y_i-\sum_{j=i+1}^n u_{ij}x_j}{u_{ii}}.
\]

这就是为什么很多算法的目标都是“把矩阵变成三角形”。

\phantomsection

\section*{一、直接法：Gauss 消元、LU、PLU}

\addcontentsline{toc}{section}{一、直接法：Gauss 消元、LU、PLU}

\phantomsection

\subsection*{1. 本章解决什么问题}

\addcontentsline{toc}{subsection}{1. 本章解决什么问题}

直接法要解决的是：

\[
Ax=b.
\]

基本思路是：用消元把 \code{\detokenize{A}} 变成上三角 \code{\detokenize{U}}，同时把消元过程记录成下三角 \code{\detokenize{L}}。最终得到：

\[
A=LU.
\]

若需要换行，则得到：

\[
PA=LU.
\]

如果还允许换列，即全主元，则写成：

\[
PAQ=LU.
\]

\phantomsection

\subsection*{2. Gauss 消元的本质}

\addcontentsline{toc}{subsection}{2. Gauss 消元的本质}

以第 \code{\detokenize{k}} 列为例，消去主元下面的元素。设当前主元是 $a_{kk}^{(k)}$，要消去第 \code{\detokenize{i}} 行第 \code{\detokenize{k}} 列元素，乘子为：

\[
l_{ik}=\frac{a_{ik}^{(k)}}{a_{kk}^{(k)}}.
\]

然后用：

\[
\text{第 } i \text{ 行}\leftarrow \text{第 } i \text{ 行}-l_{ik}\text{第 } k \text{ 行}.
\]

矩阵元素更新为：

\[
a_{ij}^{(k+1)}=a_{ij}^{(k)}-l_{ik}a_{kj}^{(k)}.
\]

这些乘子 $l_{ik}$ 放进 \code{\detokenize{L}} 的下三角部分，上三角结果就是 \code{\detokenize{U}}。

考场手算时不要只写结果。建议格式：

\begin{enumerate}

\item 写出第一步乘子；

\item 写出消元后的矩阵；

\item 写出第二步乘子；

\item 最后整理 \code{\detokenize{L}} 和 \code{\detokenize{U}}。

\end{enumerate}

\phantomsection

\subsection*{3. Doolittle 型 LU 分解}

\addcontentsline{toc}{subsection}{3. Doolittle 型 LU 分解}

Doolittle 分解约定 \code{\detokenize{L}} 是单位下三角：

\[
A=LU,\quad
L=\begin{pmatrix}
1&0&\cdots&0\\
l_{21}&1&\cdots&0\\
\vdots&\vdots&\ddots&\vdots\\
l_{n1}&l_{n2}&\cdots&1
\end{pmatrix}.
\]

如果不换行，普通 LU 存在的一个常见充分条件是所有顺序主子式非零：

\[
\det A_k\ne 0,\quad k=1,\ldots,n-1.
\]

注意这只是常见的判别说法。考试如果问“举一个非奇异但没有普通 LU 的矩阵”，可以写：

\[
A=\begin{pmatrix}0&1\\1&0\end{pmatrix}.
\]

它非奇异，但第一步主元是 0，普通 LU 不能开始。换行后可以。

\phantomsection

\subsection*{4. 用 LU 解方程}

\addcontentsline{toc}{subsection}{4. 用 LU 解方程}

如果

\[
A=LU,
\]

则

\[
Ax=b\quad \Longleftrightarrow\quad LUx=b.
\]

令

\[
Ux=y,
\]

先解

\[
Ly=b,
\]

再解

\[
Ux=y.
\]

考场模板：

\begin{enumerate}

\item $A=LU$；

\item 前代解 $Ly=b$；

\item 回代解 $Ux=y$；

\item 写出 \code{\detokenize{x}}。

\end{enumerate}

如果是列主元：

\[
PA=LU.
\]

则

\[
Ax=b\quad \Longleftrightarrow\quad PAx=Pb\quad \Longleftrightarrow\quad LUx=Pb.
\]

所以必须先把右端项也换行。常见失分点就是忘记 \code{\detokenize{Pb}}。

\phantomsection

\subsection*{5. 列主元 PLU}

\addcontentsline{toc}{subsection}{5. 列主元 PLU}

列主元 Gauss 消元每一步在当前列中，从第 \code{\detokenize{k}} 行到第 \code{\detokenize{n}} 行选绝对值最大的元素作为主元。

为什么要选主元：

\begin{itemize}

\item 避免主元为 0；

\item 避免主元太小导致乘子太大；

\item 抑制舍入误差传播；

\item 在实际计算中比普通 Gauss 消元稳定得多。

\end{itemize}

写法约定容易混淆。你可以在卷面上明确写：

\[
PA=LU.
\]

然后始终按这个约定写：

\[
LUx=Pb.
\]

如果别人写 $A=PLU$，可能是把置换矩阵放到了另一边。考试时不怕符号不同，怕前后不一致。

\phantomsection

\subsection*{6. 高频例题：普通 LU 与 PLU}

\addcontentsline{toc}{subsection}{6. 高频例题：普通 LU 与 PLU}

题目：

\[
A=\begin{pmatrix}1&4&7\\2&5&8\\3&6&10\end{pmatrix},\quad b=(1,1,1)^T.
\]

普通消元：

第一列主元为 1，乘子：

\[
l_{21}=2,\quad l_{31}=3.
\]

消元后得到：

\[
\begin{pmatrix}
1&4&7\\
0&-3&-6\\
0&-6&-11
\end{pmatrix}.
\]

第二列主元为 -3，乘子：

\[
l_{32}=2.
\]

得到：

\[
U=\begin{pmatrix}
1&4&7\\
0&-3&-6\\
0&0&1
\end{pmatrix}.
\]

所以

\[
L=\begin{pmatrix}
1&0&0\\
2&1&0\\
3&2&1
\end{pmatrix}.
\]

解方程：

\[
Ly=b.
\]

前代：

\[
y_1=1,\quad 2y_1+y_2=1,\quad 3y_1+2y_2+y_3=1,
\]

所以

\[
y=(1,-1,0)^T.
\]

再解：

\[
Ux=y.
\]

回代：

\[
x_3=0,\quad -3x_2-6x_3=-1,\quad x_1+4x_2+7x_3=1.
\]

得到：

\[
x=(-1/3,1/3,0)^T.
\]

一种列主元分解为：

\[
P=\begin{pmatrix}0&0&1\\1&0&0\\0&1&0\end{pmatrix},
\quad
L=\begin{pmatrix}1&0&0\\1/3&1&0\\2/3&1/2&1\end{pmatrix},
\quad
U=\begin{pmatrix}3&6&10\\0&2&11/3\\0&0&-1/2\end{pmatrix}.
\]

这类题的关键不是死背矩阵，而是知道每个 \code{\detokenize{L}} 中的数来自哪里。

\phantomsection

\subsection*{7. 全主元 Gauss 消元}

\addcontentsline{toc}{subsection}{7. 全主元 Gauss 消元}

全主元每一步不仅换行，也换列。第 \code{\detokenize{k}} 步在剩余子矩阵中选绝对值最大的元素作为主元。

结果写成：

\[
PAQ=LU.
\]

解方程时要小心变量顺序。若

\[
PAQ=LU,
\]

则

\[
PAx=Pb,\quad PAQ(Q^Tx)=Pb.
\]

令

\[
z=Q^Tx,
\]

先解

\[
LUz=Pb,
\]

最后还原：

\[
x=Qz.
\]

全主元比列主元更稳定，但实现更麻烦。纸面考试如果考全主元，多半是让你理解 \code{\detokenize{Q}} 代表列交换，也就是未知量顺序交换。

\phantomsection

\subsection*{8. Gauss-Jordan 求逆}

\addcontentsline{toc}{subsection}{8. Gauss-Jordan 求逆}

Gauss-Jordan 的目标是把 \code{\detokenize{A}} 化成单位阵：

\[
[A\ |\ I]\longrightarrow [I\ |\ A^{-1}].
\]

它适合手算小矩阵逆，但实际解方程不如 LU 高效。若题目问 Gauss-Jordan，答题要点是：

\begin{enumerate}

\item 对增广矩阵做初等行变换；

\item 左边化为 \code{\detokenize{I}}；

\item 右边即 $A^{-1}$；

\item 若过程中出现无法选主元，则 \code{\detokenize{A}} 奇异或需要换行。

\end{enumerate}

\phantomsection

\section*{二、三对角方程与追赶法}

\addcontentsline{toc}{section}{二、三对角方程与追赶法}

\phantomsection

\subsection*{1. 三对角矩阵}

\addcontentsline{toc}{subsection}{1. 三对角矩阵}

三对角方程形如：

\[
a_i x_{i-1}+b_i x_i+c_i x_{i+1}=d_i.
\]

矩阵形式是：

\[
\begin{pmatrix}
b_1&c_1&0&\cdots&0\\
a_2&b_2&c_2&\ddots&\vdots\\
0&a_3&b_3&\ddots&0\\
\vdots&\ddots&\ddots&\ddots&c_{n-1}\\
0&\cdots&0&a_n&b_n
\end{pmatrix}
\begin{pmatrix}x_1\\x_2\\\vdots\\x_n\end{pmatrix}
=
\begin{pmatrix}d_1\\d_2\\\vdots\\d_n\end{pmatrix}.
\]

普通 Gauss 消元会破坏大量零结构吗？对三对角矩阵不会完全破坏，但追赶法会更省事，因为它只记录三条对角线。

\phantomsection

\subsection*{2. 追赶法递推}

\addcontentsline{toc}{subsection}{2. 追赶法递推}

本质是三对角 LU 分解。设：

\[
u_1=b_1.
\]

对 $i=2,\ldots,n$：

\[
l_i=\frac{a_i}{u_{i-1}},\quad u_i=b_i-l_i c_{i-1}.
\]

前代：

\[
y_1=d_1,\quad y_i=d_i-l_i y_{i-1}.
\]

回代：

\[
x_n=\frac{y_n}{u_n},\quad
x_i=\frac{y_i-c_i x_{i+1}}{u_i}.
\]

考场如果给小三对角矩阵，直接写普通消元也能得分；但如果题目明确说追赶法，最好写出 $l_i,u_i,y_i,x_i$ 递推。

\phantomsection

\subsection*{3. 追赶法的适用提醒}

\addcontentsline{toc}{subsection}{3. 追赶法的适用提醒}

追赶法要求过程中 $u_i$ 不为零。若矩阵严格对角占优，通常可以保证算法顺利且稳定。

严格对角占优是：

\[
|a_{ii}|>\sum_{j\ne i}|a_{ij}|.
\]

它既常用于直接法可行性，也常用于迭代法收敛性。

\phantomsection

\section*{三、正定矩阵、Cholesky 与 LDLT}

\addcontentsline{toc}{section}{三、正定矩阵、Cholesky 与 LDLT}

\phantomsection

\subsection*{1. 为什么正定矩阵特殊}

\addcontentsline{toc}{subsection}{1. 为什么正定矩阵特殊}

如果 \code{\detokenize{A}} 是 SPD，那么解

\[
Ax=b
\]

可以用平方根法：

\[
A=LL^T.
\]

相比普通 LU，它利用了对称性，只需要存一半矩阵，而且稳定性好。

\phantomsection

\subsection*{2. Cholesky 分解}

\addcontentsline{toc}{subsection}{2. Cholesky 分解}

设 \code{\detokenize{A}} 是 SPD，则存在唯一的下三角矩阵 \code{\detokenize{L}}，且对角元为正，使得：

\[
A=LL^T.
\]

递推公式：

\[
l_{kk}=\sqrt{a_{kk}-\sum_{s=1}^{k-1}l_{ks}^2}.
\]

当 $i>k$ 时：

\[
l_{ik}=\frac{a_{ik}-\sum_{s=1}^{k-1}l_{is}l_{ks}}{l_{kk}}.
\]

手算步骤：

\begin{enumerate}

\item 算 $l_{11}$；

\item 用第一列公式算 $l_{21},l_{31},...$；

\item 算 $l_{22}$；

\item 继续向右下角推进；

\item 最后解 $Ly=b$ 和 $L^Tx=y$。

\end{enumerate}

\phantomsection

\subsection*{3. 改进平方根法 LDLT}

\addcontentsline{toc}{subsection}{3. 改进平方根法 LDLT}

改进平方根法写作：

\[
A=LDL^T,
\]

其中 \code{\detokenize{L}} 是单位下三角，\code{\detokenize{D}} 是对角阵。

它的优点是避免开平方。对精确手算题也常更干净，因为很多数可以留在 \code{\detokenize{D}} 中。

递推形式：

\[
d_k=a_{kk}-\sum_{s=1}^{k-1}l_{ks}^2d_s.
\]

当 $i>k$ 时：

\[
l_{ik}=\frac{a_{ik}-\sum_{s=1}^{k-1}l_{is}l_{ks}d_s}{d_k}.
\]

解方程：

\[
LDL^Tx=b.
\]

令

\[
Ly=b,\quad Dz=y,\quad L^Tx=z.
\]

三步都简单：前代、对角缩放、回代。

\phantomsection

\subsection*{4. 正定性证明模板}

\addcontentsline{toc}{subsection}{4. 正定性证明模板}

题目让“证明 \code{\detokenize{A}} 正定”，不要上来就分解。先证明正定，再用分解。

模板一：顺序主子式。

\begin{enumerate}

\item 先说 $A=A^T$；

\item 算顺序主子式；

\item 全部大于零；

\item 所以 \code{\detokenize{A}} SPD。

\end{enumerate}

模板二：特征值。

\begin{enumerate}

\item 算出或说明所有特征值；

\item 全部大于零；

\item 对称矩阵特征值全正等价于 SPD。

\end{enumerate}

模板三：二次型配方。

写：

\[
x^TAx=\text{若干平方和},
\]

并说明只要 $x\ne 0$，平方和大于 0。

\phantomsection

\subsection*{5. 高频例题：Cholesky}

\addcontentsline{toc}{subsection}{5. 高频例题：Cholesky}

题目：

\[
A=\begin{pmatrix}
4&-2&4&2\\
-2&10&-2&-7\\
4&-2&8&4\\
2&-7&4&7
\end{pmatrix},\quad b=(8,2,16,6)^T.
\]

一个 Cholesky 分解是：

\[
L=\begin{pmatrix}
2&0&0&0\\
-1&3&0&0\\
2&0&2&0\\
1&-2&1&1
\end{pmatrix}.
\]

可以检查：

\[
LL^T=A.
\]

解方程时：

\[
Ly=b,\quad L^Tx=y.
\]

最终：

\[
x=(1,2,1,2)^T.
\]

这题的考试价值在于：先判正定，再分解，再前代回代，完整链条不能断。

\phantomsection

\subsection*{6. 常见失分点}

\addcontentsline{toc}{subsection}{6. 常见失分点}

\begin{itemize}

\item 没证明对称就直接说正定；

\item Cholesky 用在非正定矩阵上；

\item $LL^T$ 写成 $L^TL$；

\item $LDL^T$ 中忘记 \code{\detokenize{L}} 是单位下三角；

\item 解方程时只分解，不做回代。

\end{itemize}

\phantomsection

\section*{四、向量范数、矩阵范数与条件数}

\addcontentsline{toc}{section}{四、向量范数、矩阵范数与条件数}

\phantomsection

\subsection*{1. 为什么要学范数}

\addcontentsline{toc}{subsection}{1. 为什么要学范数}

范数是“大小”的抽象。数值计算里我们关心：

\begin{itemize}

\item 向量误差有多大；

\item 矩阵扰动有多大；

\item 残量小是否代表误差小；

\item 算法是否会放大误差。

\end{itemize}

所以范数是误差分析和收敛分析的共同语言。

\phantomsection

\subsection*{2. 向量范数定义}

\addcontentsline{toc}{subsection}{2. 向量范数定义}

一个函数 $\|x\|$ 是向量范数，需要满足三条：

\begin{enumerate}

\item 非负性：$\|x\|>=0$，且 $\|x\|=0$ 当且仅当 $x=0$；

\item 齐次性：$\|\alpha x\|=|\alpha| \|x\|$；

\item 三角不等式：$\|x+y\|<=\|x\|+\|y\|$。

\end{enumerate}

常见范数：

\[
\|x\|_1=\sum_i |x_i|.
\]

\[
\|x\|_2=\left(\sum_i |x_i|^2\right)^{1/2}.
\]

\[
\|x\|_\infty=\max_i |x_i|.
\]

一般 \code{\detokenize{p}} 范数：

\[
\|x\|_p=\left(\sum_i |x_i|^p\right)^{1/p},\quad p\ge 1.
\]

常用等价关系：

\[
\|x\|_\infty \le \|x\|_2 \le \sqrt n\|x\|_\infty.
\]

\[
\|x\|_2 \le \|x\|_1 \le \sqrt n\|x\|_2.
\]

\[
\|x\|_\infty \le \|x\|_1 \le n\|x\|_\infty.
\]

这些不等式常用于证明收敛或估计误差。

\phantomsection

\subsection*{3. 矩阵范数定义}

\addcontentsline{toc}{subsection}{3. 矩阵范数定义}

矩阵范数至少要满足：

\begin{enumerate}

\item 非负性；

\item 齐次性；

\item 三角不等式；

\item 次乘性：

\end{enumerate}

\[
\|AB\|\le \|A\|\|B\|.
\]

如果是由向量范数诱导出的矩阵范数：

\[
\|A\|=\max_{x\ne 0}\frac{\|Ax\|}{\|x\|}.
\]

则自动满足：

\[
\|Ax\|\le \|A\|\|x\|.
\]

常见诱导矩阵范数：

\[
\|A\|_1=\max_j \sum_i |a_{ij}|.
\]

也就是最大列和。

\[
\|A\|_\infty=\max_i \sum_j |a_{ij}|.
\]

也就是最大行和。

谱范数：

\[
\|A\|_2=\sqrt{\lambda_{\max}(A^TA)}.
\]

如果 \code{\detokenize{A}} 对称，则：

\[
\|A\|_2=\max_i |\lambda_i(A)|.
\]

\phantomsection

\subsection*{4. Frobenius 范数}

\addcontentsline{toc}{subsection}{4. Frobenius 范数}

Frobenius 范数定义为：

\[
\|A\|_F=\left(\sum_{i,j}|a_{ij}|^2\right)^{1/2}.
\]

也可写成：

\[
\|A\|_F=\sqrt{\operatorname{tr}(A^TA)}.
\]

证明它是矩阵范数，最关键的是次乘性。常用方法是把列向量拆开，或使用 Cauchy-Schwarz 不等式。考试更常见的是证明正交不变性：

\[
\|U^TAV\|_F^2
=\operatorname{tr}(V^TA^TUU^TAV)
=\operatorname{tr}(V^TA^TAV).
\]

若 \code{\detokenize{V}} 正交，则利用迹的循环性质：

\[
\operatorname{tr}(V^TA^TAV)=\operatorname{tr}(A^TAVV^T)=\operatorname{tr}(A^TA).
\]

所以：

\[
\|U^TAV\|_F=\|A\|_F.
\]

\phantomsection

\subsection*{5. 相似诱导范数}

\addcontentsline{toc}{subsection}{5. 相似诱导范数}

历年高频题：

\[
\|A\|_X=\|X^{-1}AX\|_2.
\]

其中 \code{\detokenize{X}} 可逆。证明它是矩阵范数。

答题模板：

非负性来自 2-范数：

\[
\|A\|_X=\|X^{-1}AX\|_2\ge 0.
\]

若 $\|A\|_X=0$，则

\[
X^{-1}AX=0.
\]

左右乘 \code{\detokenize{X}} 和 $X^{-1}$ 得：

\[
A=0.
\]

齐次性：

\[
\|\alpha A\|_X=\|X^{-1}(\alpha A)X\|_2=|\alpha|\|X^{-1}AX\|_2.
\]

三角不等式：

\[
\|A+B\|_X=\|X^{-1}(A+B)X\|_2
\le \|X^{-1}AX\|_2+\|X^{-1}BX\|_2.
\]

次乘性：

\[
\|AB\|_X=\|X^{-1}ABX\|_2
=\|(X^{-1}AX)(X^{-1}BX)\|_2
\le \|A\|_X\|B\|_X.
\]

这题的核心是中间插入：

\[
XX^{-1}=I.
\]

\phantomsection

\subsection*{6. 谱半径与收敛矩阵}

\addcontentsline{toc}{subsection}{6. 谱半径与收敛矩阵}

谱半径定义为：

\[
\rho(A)=\max_i |\lambda_i(A)|.
\]

结论：

\[
A^k\to 0\quad \Longleftrightarrow\quad \rho(A)<1.
\]

这就是迭代法的核心。只要能把误差写成：

\[
e^{(k+1)}=Be^{(k)},
\]

就有：

\[
e^{(k)}=B^k e^{(0)}.
\]

所以收敛充要条件是：

\[
\rho(B)<1.
\]

\phantomsection

\subsection*{7. 条件数}

\addcontentsline{toc}{subsection}{7. 条件数}

条件数衡量问题本身对扰动的敏感程度：

\[
\kappa(A)=\|A\|\|A^{-1}\|.
\]

若 \code{\detokenize{A}} 非奇异，且只扰动右端项：

\[
A(x+\delta x)=b+\delta b.
\]

有粗略误差界：

\[
\frac{\|\delta x\|}{\|x\|}
\lesssim
\kappa(A)\frac{\|\delta b\|}{\|b\|}.
\]

如果条件数很大，哪怕残量小，误差也可能大。

残量是：

\[
r=b-A\hat x.
\]

误差是：

\[
e=x-\hat x.
\]

二者关系：

\[
Ae=r.
\]

所以：

\[
e=A^{-1}r.
\]

于是：

\[
\|e\|\le \|A^{-1}\|\|r\|.
\]

这解释了为什么病态矩阵中“残量小不等于误差小”。

\phantomsection

\subsection*{8. 舍入误差与增长因子}

\addcontentsline{toc}{subsection}{8. 舍入误差与增长因子}

浮点运算会产生舍入误差。数值算法的好坏，很大程度上看它会不会放大这些误差。

Gauss 消元中如果主元很小，乘子：

\[
l_{ik}=\frac{a_{ik}^{(k)}}{a_{kk}^{(k)}}
\]

可能很大，导致舍入误差被放大。选主元就是为了控制乘子大小。

增长因子可以粗略理解为消元过程中元素最大值相对原矩阵最大值的放大倍数。它越大，潜在误差放大越严重。

考场遇到“为什么列主元稳定”，不要写玄学句子，写：

\begin{itemize}

\item 选择绝对值较大的主元；

\item 控制乘子；

\item 减少舍入误差放大；

\item 避免零主元。

\end{itemize}

\phantomsection

\section*{五、最小二乘、正则方程与投影}

\addcontentsline{toc}{section}{五、最小二乘、正则方程与投影}

\phantomsection

\subsection*{1. 最小二乘问题是什么}

\addcontentsline{toc}{subsection}{1. 最小二乘问题是什么}

当方程组 $Ax=b$ 没有精确解时，想找一个 \code{\detokenize{x}} 让残差最小：

\[
\min_x \|Ax-b\|_2.
\]

这里 \code{\detokenize{A}} 通常是高矩阵：

\[
A\in\mathbb R^{m\times n},\quad m\ge n.
\]

如果 \code{\detokenize{A}} 满列秩，则最小二乘解唯一。

几何理解：\code{\detokenize{Ax}} 永远在列空间 \code{\detokenize{R(A)}} 里。我们要找列空间里离 \code{\detokenize{b}} 最近的向量。

最近点是 \code{\detokenize{b}} 在 \code{\detokenize{R(A)}} 上的正交投影。

\phantomsection

\subsection*{2. 正则方程}

\addcontentsline{toc}{subsection}{2. 正则方程}

最优残差：

\[
r=b-Ax.
\]

最优条件是残差垂直于列空间：

\[
A^Tr=0.
\]

即：

\[
A^T(b-Ax)=0.
\]

所以：

\[
A^TAx=A^Tb.
\]

这就是正则方程。

如果 \code{\detokenize{A}} 满列秩，则 $A^TA$ SPD。证明：

对任意 $x\ne 0$，

\[
x^TA^TAx=\|Ax\|_2^2.
\]

满列秩说明 $Ax\ne 0$，因此：

\[
\|Ax\|_2^2>0.
\]

所以 $A^TA$ 正定。

\phantomsection

\subsection*{3. 正则方程的优缺点}

\addcontentsline{toc}{subsection}{3. 正则方程的优缺点}

优点：

\begin{itemize}

\item 推导简单；

\item 变成方阵 SPD 系统；

\item 可以用 Cholesky 解。

\end{itemize}

缺点：

\begin{itemize}

\item 条件数会平方；

\end{itemize}

\[
\kappa_2(A^TA)=\kappa_2(A)^2.
\]

所以数值稳定性不如 QR。

考试如果问“为什么 QR 更好”，可以答：正交变换不改变 2-范数，不会像正则方程那样平方条件数。

\phantomsection

\subsection*{4. 正交投影矩阵}

\addcontentsline{toc}{subsection}{4. 正交投影矩阵}

若 \code{\detokenize{A}} 满列秩，则到列空间 \code{\detokenize{R(A)}} 的正交投影矩阵是：

\[
P=A(A^TA)^{-1}A^T.
\]

它满足：

\[
P^T=P,\quad P^2=P.
\]

证明对称：

\[
P^T=(A(A^TA)^{-1}A^T)^T
=A((A^TA)^{-1})^TA^T=P.
\]

因为 $A^TA$ 对称，其逆也对称。

证明幂等：

\[
P^2=A(A^TA)^{-1}A^TA(A^TA)^{-1}A^T=P.
\]

如果 \code{\detokenize{P}} 是非零正交投影，则特征值只可能是 0 或 1，因此：

\[
\|P\|_2=1.
\]

历年题常见问法：证明

\[
\|A(A^TA)^{-1}A^T\|_2=1.
\]

答题路线就是：先证明它是正交投影，再说谱范数等于最大特征值绝对值。

\phantomsection

\subsection*{5. 高频例题：正则方程与平方根法}

\addcontentsline{toc}{subsection}{5. 高频例题：正则方程与平方根法}

题目：

\[
A=\begin{pmatrix}0&-1\\3&2\\0&-2\end{pmatrix},\quad
b=\begin{pmatrix}-1/3\\5/3\\-2/3\end{pmatrix}.
\]

计算：

\[
A^TA=
\begin{pmatrix}
9&6\\
6&9
\end{pmatrix}.
\]

再算：

\[
A^Tb=
\begin{pmatrix}5\\5\end{pmatrix}.
\]

正则方程：

\[
\begin{pmatrix}9&6\\6&9\end{pmatrix}x=
\begin{pmatrix}5\\5\end{pmatrix}.
\]

解得：

\[
x=(1/3,1/3)^T.
\]

如果题目要求平方根法，可以对 $A^TA$ 做 Cholesky。因为

\[
\begin{pmatrix}9&6\\6&9\end{pmatrix}
=
\begin{pmatrix}3&0\\2&\sqrt5\end{pmatrix}
\begin{pmatrix}3&2\\0&\sqrt5\end{pmatrix}.
\]

然后前代回代即可。

\phantomsection

\subsection*{6. 广义逆与最小二乘}

\addcontentsline{toc}{subsection}{6. 广义逆与最小二乘}

若 \code{\detokenize{A}} 满列秩，Moore-Penrose 逆为：

\[
A^+=(A^TA)^{-1}A^T.
\]

最小二乘解可以写成：

\[
x=A^+b.
\]

如果题目说“存在矩阵 \code{\detokenize{X}}，使得对任意 \code{\detokenize{b}}，$x=Xb$ 都是最小二乘解”，那么 \code{\detokenize{AX}} 是投影矩阵。常见结论：

\[
AXA=A,\quad (AX)^T=AX.
\]

理解：\code{\detokenize{AXb}} 是 \code{\detokenize{b}} 在 \code{\detokenize{R(A)}} 上的正交投影，所以 \code{\detokenize{AX}} 对称且幂等，并且在 \code{\detokenize{R(A)}} 上保持不变。

\phantomsection

\section*{六、QR 分解与正交化方法}

\addcontentsline{toc}{section}{六、QR 分解与正交化方法}

\phantomsection

\subsection*{1. QR 为什么适合最小二乘}

\addcontentsline{toc}{subsection}{1. QR 为什么适合最小二乘}

如果 \code{\detokenize{A}} 满列秩，可以分解：

\[
A=QR,
\]

其中 \code{\detokenize{Q}} 列正交，\code{\detokenize{R}} 上三角。

更完整地写：

\[
A=Q\begin{pmatrix}R\\0\end{pmatrix},\quad Q=[Q_1,Q_2].
\]

因为 \code{\detokenize{Q}} 正交：

\[
\|Ax-b\|_2=\|Q^T(Ax-b)\|_2.
\]

于是：

\[
\|Ax-b\|_2^2
=\left\|\begin{pmatrix}R\\0\end{pmatrix}x-
\begin{pmatrix}Q_1^Tb\\Q_2^Tb\end{pmatrix}\right\|_2^2.
\]

也就是：

\[
\|Ax-b\|_2^2
=\|Rx-Q_1^Tb\|_2^2+\|Q_2^Tb\|_2^2.
\]

第二项与 \code{\detokenize{x}} 无关，所以最小化第一项：

\[
Rx=Q_1^Tb.
\]

这就是 QR 解最小二乘的核心。

\phantomsection

\subsection*{2. Gram-Schmidt 正交化}

\addcontentsline{toc}{subsection}{2. Gram-Schmidt 正交化}

给定线性无关向量 $a_1,\ldots,a_n$，想构造正交单位向量 $q_1,\ldots,q_n$。

经典 Gram-Schmidt：

\[
r_{11}=\|a_1\|_2,\quad q_1=\frac{a_1}{r_{11}}.
\]

对第 \code{\detokenize{k}} 个向量：

\[
r_{ik}=q_i^Ta_k,\quad i=1,\ldots,k-1.
\]

去掉已有方向投影：

\[
v_k=a_k-\sum_{i=1}^{k-1}r_{ik}q_i.
\]

归一化：

\[
r_{kk}=\|v_k\|_2,\quad q_k=\frac{v_k}{r_{kk}}.
\]

最后有：

\[
A=QR.
\]

经典 Gram-Schmidt 容易因舍入误差损失正交性；改进 Gram-Schmidt 数值上更好。纸面考试通常考公式和思想。

\phantomsection

\subsection*{3. Householder 变换}

\addcontentsline{toc}{subsection}{3. Householder 变换}

Householder 变换定义：

\[
H=I-2ww^T,\quad \|w\|_2=1.
\]

性质：

\[
H^T=H,\quad H^TH=I,\quad H^2=I.
\]

它是关于某个超平面的反射。它的用途是“一次消掉一串元素”。

给向量 \code{\detokenize{x}}，想构造 \code{\detokenize{H}} 使：

\[
Hx=\alpha e_1.
\]

取：

\[
\alpha=\pm \|x\|_2,
\quad
w=\frac{x-\alpha e_1}{\|x-\alpha e_1\|_2}.
\]

实际计算中选择 $\alpha$ 的符号要避免相消。常用建议：

\[
\alpha=-\operatorname{sign}(x_1)\|x\|_2.
\]

考试中如果只要构造，可以选方便符号，但要保证结果对。

\phantomsection

\subsection*{4. Householder 做 QR}

\addcontentsline{toc}{subsection}{4. Householder 做 QR}

对矩阵 \code{\detokenize{A}}，第 1 步取第一列下方向量，构造 Householder 把第一列第 2 到第 \code{\detokenize{m}} 个元素消掉。

第 2 步在右下角子矩阵中继续，只动第 2 列下面部分。

最后：

\[
H_k\cdots H_2H_1A=R.
\]

因为 Householder 正交，令：

\[
Q=H_1H_2\cdots H_k.
\]

得到：

\[
A=QR.
\]

\phantomsection

\subsection*{5. 高频 Householder 例题}

\addcontentsline{toc}{subsection}{5. 高频 Householder 例题}

题目中常出现：

\[
x=(3,1,6,4,2,2)^T.
\]

要求构造正交变换使：

\[
Hx=(3,\alpha,4,6,0,0)^T.
\]

观察：第 1、3、4 个分量不变，只作用在第 2、5、6 个分量上。对应子向量是：

\[
y=(1,2,2)^T.
\]

希望：

\[
H_3y=(3,0,0)^T.
\]

可取：

\[
H_3=
\begin{pmatrix}
1/3&2/3&2/3\\
2/3&1/3&-2/3\\
2/3&-2/3&1/3
\end{pmatrix}.
\]

然后把 $H_3$ 嵌入到第 2、5、6 个坐标对应的位置即可。

\phantomsection

\subsection*{6. Givens 旋转}

\addcontentsline{toc}{subsection}{6. Givens 旋转}

Givens 旋转是二维旋转：

\[
G=
\begin{pmatrix}
c&s\\
-s&c
\end{pmatrix},\quad c^2+s^2=1.
\]

它的用途是“一次消掉一个元素”。

若要把：

\[
\begin{pmatrix}a\\b\end{pmatrix}
\]

化为：

\[
\begin{pmatrix}r\\0\end{pmatrix},
\]

可取：

\[
r=\sqrt{a^2+b^2},\quad c=\frac a r,\quad s=\frac b r.
\]

注意符号约定和你写的 \code{\detokenize{G}} 有关。只要最终乘出来正确即可。

\phantomsection

\subsection*{7. 高频 Givens 例题}

\addcontentsline{toc}{subsection}{7. 高频 Givens 例题}

题目：

\[
\begin{pmatrix}c&s\\-s&c\end{pmatrix}
\begin{pmatrix}7\\-1\end{pmatrix}
=
\begin{pmatrix}\beta\\\beta\end{pmatrix}.
\]

列方程：

\[
7c-s=\beta,
\]

\[
-7s-c=\beta,
\]

\[
c^2+s^2=1.
\]

一组解是：

\[
c=3/5,\quad s=-4/5,\quad \beta=5.
\]

考试中如果给等价符号解，也要检查乘法是否真的得到两个相同分量。

\phantomsection

\subsection*{8. Householder 与 Givens 对比}

\addcontentsline{toc}{subsection}{8. Householder 与 Givens 对比}

Householder：

\begin{itemize}

\item 一次消掉一串；

\item 适合稠密矩阵；

\item 常用于 QR 和 Hessenberg 化。

\end{itemize}

Givens：

\begin{itemize}

\item 一次消掉一个；

\item 适合稀疏矩阵或只需要局部修改；

\item 常用于 QR 更新和手算旋转题。

\end{itemize}

共同点：

\begin{itemize}

\item 都是正交变换；

\item 都保持 2-范数；

\item 都有助于数值稳定。

\end{itemize}

\phantomsection

\section*{七、定常迭代法：Jacobi、Gauss-Seidel、SOR}

\addcontentsline{toc}{section}{七、定常迭代法：Jacobi、Gauss-Seidel、SOR}

\phantomsection

\subsection*{1. 迭代法的统一思想}

\addcontentsline{toc}{subsection}{1. 迭代法的统一思想}

把线性方程组：

\[
Ax=b
\]

改写成：

\[
x=Bx+g.
\]

迭代格式：

\[
x^{(k+1)}=Bx^{(k)}+g.
\]

设精确解满足：

\[
x=Bx+g.
\]

误差：

\[
e^{(k)}=x^{(k)}-x.
\]

相减得到：

\[
e^{(k+1)}=Be^{(k)}.
\]

所以：

\[
e^{(k)}=B^ke^{(0)}.
\]

收敛充要条件：

\[
\rho(B)<1.
\]

这就是所有定常迭代题的中心。

\phantomsection

\subsection*{2. 矩阵分裂约定}

\addcontentsline{toc}{subsection}{2. 矩阵分裂约定}

本课程常用约定：

\[
A=D-L-U.
\]

其中 \code{\detokenize{D}} 是对角部分，\code{\detokenize{-L}} 是严格下三角部分，\code{\detokenize{-U}} 是严格上三角部分。

也就是说，如果原矩阵非对角元素是 $a_{ij}$，那么 \code{\detokenize{L}} 和 \code{\detokenize{U}} 中往往带负号。不同教材可能写 $A=D+L+U$，所以一定要跟题目或讲义约定一致。

\phantomsection

\subsection*{3. Jacobi 迭代}

\addcontentsline{toc}{subsection}{3. Jacobi 迭代}

从

\[
(D-L-U)x=b
\]

得：

\[
Dx=(L+U)x+b.
\]

迭代：

\[
x^{(k+1)}=D^{-1}(L+U)x^{(k)}+D^{-1}b.
\]

迭代矩阵：

\[
B_J=D^{-1}(L+U).
\]

如果用原矩阵非对角部分写，也常见：

\[
B_J=-D^{-1}(A-D).
\]

手算时最稳的方法是逐行写：

\[
x_i^{(k+1)}=\frac{1}{a_{ii}}\left(b_i-\sum_{j\ne i}a_{ij}x_j^{(k)}\right).
\]

Jacobi 的特点是每个分量都用上一轮值。

\phantomsection

\subsection*{4. Gauss-Seidel 迭代}

\addcontentsline{toc}{subsection}{4. Gauss-Seidel 迭代}

从

\[
(D-L)x=Ux+b
\]

得：

\[
x^{(k+1)}=(D-L)^{-1}Ux^{(k)}+(D-L)^{-1}b.
\]

迭代矩阵：

\[
B_{GS}=(D-L)^{-1}U.
\]

逐行理解：第 \code{\detokenize{i}} 个分量更新时，前面已经更新过的分量用新值，后面还没更新的分量用旧值。

\phantomsection

\subsection*{5. SOR 迭代}

\addcontentsline{toc}{subsection}{5. SOR 迭代}

SOR 是 Gauss-Seidel 的松弛版本。参数 $\omega$ 控制更新步长。

矩阵形式：

\[
x^{(k+1)}
=(D-\omega L)^{-1}[(1-\omega)D+\omega U]x^{(k)}
+\omega(D-\omega L)^{-1}b.
\]

迭代矩阵：

\[
B_\omega=(D-\omega L)^{-1}[(1-\omega)D+\omega U].
\]

当：

\[
\omega=1,
\]

SOR 退化为 Gauss-Seidel。

若 \code{\detokenize{A}} SPD，SOR 收敛的经典条件是：

\[
0<\omega<2.
\]

\phantomsection

\subsection*{6. 常用收敛充分条件}

\addcontentsline{toc}{subsection}{6. 常用收敛充分条件}

严格对角占优：

\[
|a_{ii}|>\sum_{j\ne i}|a_{ij}|.
\]

若 \code{\detokenize{A}} 严格对角占优，则 Jacobi 和 Gauss-Seidel 通常收敛。

若 \code{\detokenize{A}} SPD，则 Gauss-Seidel 收敛。

若 \code{\detokenize{A}} SPD 且：

\[
0<\omega<2,
\]

则 SOR 收敛。

但参数题最稳的做法仍然是：写出迭代矩阵，求特征值，判断谱半径。

\phantomsection

\subsection*{7. 参数矩阵高频模板}

\addcontentsline{toc}{subsection}{7. 参数矩阵高频模板}

题目：

\[
A=\begin{pmatrix}1&0&a\\0&1&0\\a&0&1\end{pmatrix}.
\]

正定性：

顺序主子式：

\[
1>0,\quad
\det\begin{pmatrix}1&0\\0&1\end{pmatrix}=1>0,
\]

\[
\det A=1-a^2.
\]

所以正定当且仅当：

\[
|a|<1.
\]

Jacobi 迭代矩阵：

\[
B_J=
\begin{pmatrix}
0&0&-a\\
0&0&0\\
-a&0&0
\end{pmatrix}.
\]

它的特征值是：

\[
0,\quad a,\quad -a.
\]

所以 Jacobi 收敛当且仅当：

\[
|a|<1.
\]

Gauss-Seidel 迭代矩阵的特征值为：

\[
0,\quad 0,\quad a^2.
\]

所以收敛当且仅当：

\[
|a|<1.
\]

这道题的特点是：正定、Jacobi、G-S 三问答案一样，但理由不同。

\phantomsection

\subsection*{8. 三阶矩阵写 Jacobi 迭代矩阵}

\addcontentsline{toc}{subsection}{8. 三阶矩阵写 Jacobi 迭代矩阵}

给：

\[
A=(a_{ij})_{3\times 3}.
\]

Jacobi 迭代矩阵可以直接写：

\[
B_J=
\begin{pmatrix}
0&-a_{12}/a_{11}&-a_{13}/a_{11}\\
-a_{21}/a_{22}&0&-a_{23}/a_{22}\\
-a_{31}/a_{33}&-a_{32}/a_{33}&0
\end{pmatrix}.
\]

然后求特征值或谱半径。

历年常见结论：

\[
\rho(B_1)=\sqrt5/2>1,
\quad
\rho(B_2)=0.
\]

一个谱半径大于 1，不收敛；一个谱半径等于 0，有限步内快速稳定。

\phantomsection

\subsection*{9. 误差估计}

\addcontentsline{toc}{subsection}{9. 误差估计}

如果存在某个矩阵范数使：

\[
\|B\|<1,
\]

则：

\[
\|e^{(k)}\|\le \|B\|^k\|e^{(0)}\|.
\]

相邻迭代差也可估计误差：

\[
\|x-x^{(k)}\|\le \frac{\|B\|}{1-\|B\|}\|x^{(k)}-x^{(k-1)}\|.
\]

考试不一定要求背完整公式，但要会说：迭代误差由 $B^k$ 控制，谱半径决定最终收敛。

\phantomsection

\section*{八、共轭梯度法 CG}

\addcontentsline{toc}{section}{八、共轭梯度法 CG}

\phantomsection

\subsection*{1. CG 解决什么问题}

\addcontentsline{toc}{subsection}{1. CG 解决什么问题}

CG 用于解 SPD 方程组：

\[
Ax=b,\quad A=A^T>0.
\]

它等价于最小化二次函数：

\[
\phi(x)=\frac12 x^TAx-b^Tx.
\]

梯度为：

\[
\nabla\phi(x)=Ax-b.
\]

残量定义：

\[
r_k=b-Ax_k.
\]

所以残量等于负梯度：

\[
r_k=-\nabla\phi(x_k).
\]

\phantomsection

\subsection*{2. A-内积和 A-共轭}

\addcontentsline{toc}{subsection}{2. A-内积和 A-共轭}

因为 \code{\detokenize{A}} SPD，可以定义 A-内积：

\[
(x,y)_A=x^TAy.
\]

对应 A-范数：

\[
\|x\|_A=\sqrt{x^TAx}.
\]

向量 $p_i,p_j$ A-共轭是指：

\[
p_i^TAp_j=0,\quad i\ne j.
\]

CG 的关键是选择一组 A-共轭方向。沿这些方向逐步最小化，理论上 \code{\detokenize{n}} 步内得到精确解。

\phantomsection

\subsection*{3. A-共轭向量线性无关证明}

\addcontentsline{toc}{subsection}{3. A-共轭向量线性无关证明}

命题：若 $p_1,\ldots,p_m$ 非零且两两 A-共轭，则它们线性无关。

证明：

假设：

\[
\sum_{i=1}^m c_i p_i=0.
\]

左乘：

\[
p_j^TA.
\]

得到：

\[
\sum_{i=1}^m c_i p_j^TAp_i=0.
\]

由于 A-共轭，除 $i=j$ 外全部为 0：

\[
c_j p_j^TAp_j=0.
\]

因为 \code{\detokenize{A}} SPD 且 $p_j\ne 0$：

\[
p_j^TAp_j>0.
\]

所以：

\[
c_j=0.
\]

任意 \code{\detokenize{j}} 都成立，故线性无关。

\phantomsection

\subsection*{4. 逆矩阵展开式}

\addcontentsline{toc}{subsection}{4. 逆矩阵展开式}

若 $p_1,\ldots,p_n$ 是一组 A-共轭基，则：

\[
A^{-1}=\sum_{k=1}^n\frac{p_kp_k^T}{p_k^TAp_k}.
\]

证明思路：

任意向量 \code{\detokenize{x}} 可以展开为：

\[
x=\sum_{k=1}^n \alpha_k p_k.
\]

左乘 $p_j^TA$：

\[
p_j^TAx=\alpha_j p_j^TAp_j.
\]

所以：

\[
\alpha_j=\frac{p_j^TAx}{p_j^TAp_j}.
\]

因此：

\[
x=\sum_{k=1}^n p_k\frac{p_k^TAx}{p_k^TAp_k}.
\]

令：

\[
x=A^{-1}y.
\]

则：

\[
A^{-1}y=\sum_{k=1}^n p_k\frac{p_k^Ty}{p_k^TAp_k}.
\]

对任意 \code{\detokenize{y}} 成立，所以：

\[
A^{-1}=\sum_{k=1}^n\frac{p_kp_k^T}{p_k^TAp_k}.
\]

这是历年卷中非常典型的证明题。

\phantomsection

\subsection*{5. CG 算法递推}

\addcontentsline{toc}{subsection}{5. CG 算法递推}

初始化：

\[
r_0=b-Ax_0,\quad p_0=r_0.
\]

步长：

\[
\alpha_k=\frac{r_k^Tr_k}{p_k^TAp_k}.
\]

更新解：

\[
x_{k+1}=x_k+\alpha_k p_k.
\]

更新残量：

\[
r_{k+1}=r_k-\alpha_k Ap_k.
\]

更新方向：

\[
\beta_k=\frac{r_{k+1}^Tr_{k+1}}{r_k^Tr_k}.
\]

\[
p_{k+1}=r_{k+1}+\beta_k p_k.
\]

考试手算 CG 时，建议每一步列出：

\begin{enumerate}

\item $r_k$；

\item $p_k$；

\item $Ap_k$；

\item $alpha_k$；

\item $x_{k+1}$；

\item $r_{k+1}$；

\item $beta_k$；

\item $p_{k+1}$。

\end{enumerate}

\phantomsection

\subsection*{6. CG 的正交关系}

\addcontentsline{toc}{subsection}{6. CG 的正交关系}

CG 中有两类重要关系：

\begin{itemize}

\item 残量正交：

\end{itemize}

\[
r_i^Tr_j=0,\quad i\ne j.
\]

\begin{itemize}

\item 搜索方向 A-共轭：

\end{itemize}

\[
p_i^TAp_j=0,\quad i\ne j.
\]

这解释了为什么 CG 在精确算术下最多 \code{\detokenize{n}} 步收敛。

讲义中常见关系：

\[
(r^{(k)},r^{(k)})=-\alpha_{k-1}(r^{(k)},Ap^{(k-1)}),
\]

\[
(r^{(k)},r^{(k)})=\alpha_k(p^{(k)},Ap^{(k)}).
\]

证明这类式子时，通常从：

\[
r_{k+1}=r_k-\alpha_kAp_k
\]

和：

\[
p_k=r_k+\beta_{k-1}p_{k-1}
\]

出发，再代入残量正交、方向共轭。

\phantomsection

\subsection*{7. 最速下降与 CG 的区别}

\addcontentsline{toc}{subsection}{7. 最速下降与 CG 的区别}

最速下降每步沿负梯度方向，也就是残量方向：

\[
p_k=r_k.
\]

CG 会修正方向：

\[
p_{k+1}=r_{k+1}+\beta_kp_k,
\]

使方向保持 A-共轭。

因此 CG 通常比最速下降快得多。最速下降可能“之”字形震荡，CG 用共轭方向避免重复搜索。

\phantomsection

\section*{九、特征值问题：幂法、Hessenberg、QR、Schur}

\addcontentsline{toc}{section}{九、特征值问题：幂法、Hessenberg、QR、Schur}

\phantomsection

\subsection*{1. 特征值问题的目标}

\addcontentsline{toc}{subsection}{1. 特征值问题的目标}

特征值问题是：

\[
Ax=\lambda x.
\]

直接求根不稳定也不经济。数值算法的思路是通过相似变换把矩阵化成更简单的形状。

相似变换：

\[
B=S^{-1}AS.
\]

\code{\detokenize{A}} 和 \code{\detokenize{B}} 有相同特征值，因为：

\[
\det(\lambda I-B)=\det(S^{-1}(\lambda I-A)S)=\det(\lambda I-A).
\]

若 \code{\detokenize{S}} 是正交矩阵，变换更稳定：

\[
B=Q^TAQ.
\]

\phantomsection

\subsection*{2. 幂法}

\addcontentsline{toc}{subsection}{2. 幂法}

幂法用于求模最大的特征值。假设：

\[
|\lambda_1|>|\lambda_2|\ge \cdots.
\]

若初始向量在第一特征向量方向上有非零分量，则：

\[
A^kx_0
=c_1\lambda_1^kv_1+c_2\lambda_2^kv_2+\cdots.
\]

随着 \code{\detokenize{k}} 增大，第一项占主导。

实际算法每步要归一化，避免数值溢出。

幂法可能失败或表现异常的情况：

\begin{itemize}

\item 最大模特征值不唯一；

\item 矩阵不可对角化且有 Jordan 块；

\item 初始向量恰好没有主特征向量分量；

\item 最大特征值有相同模但符号相反，例如 $\lambda$ 和 $-\lambda$。

\end{itemize}

\phantomsection

\subsection*{3. 逆迭代与位移}

\addcontentsline{toc}{subsection}{3. 逆迭代与位移}

若想找靠近某个 $\mu$ 的特征值，可以对：

\[
(A-\mu I)^{-1}
\]

做幂法。

因为若 \code{\detokenize{A}} 的特征值是 $lambda_i$，则 $(A-\\mu I)^{-1}$ 的特征值是：

\[
\frac{1}{\lambda_i-\mu}.
\]

离 $\mu$ 最近的特征值会变成模最大的特征值。

实际每步解线性方程：

\[
(A-\mu I)y_k=x_k.
\]

这叫逆迭代。若 $\mu$ 动态取 Rayleigh 商，则可得到 Rayleigh 商迭代。

\phantomsection

\subsection*{4. Householder 上 Hessenberg 化}

\addcontentsline{toc}{subsection}{4. Householder 上 Hessenberg 化}

上 Hessenberg 矩阵定义：

\[
h_{ij}=0,\quad i>j+1.
\]

也就是主对角线下一条次对角线以下全为 0。

一般矩阵可以用正交相似变换化为上 Hessenberg：

\[
Q^TAQ=H.
\]

第 \code{\detokenize{k}} 步：

\begin{enumerate}

\item 取第 \code{\detokenize{k}} 列从第 \code{\detokenize{k+1}} 行到第 \code{\detokenize{n}} 行的向量；

\item 构造 Householder，把该向量除第一个分量外消为 0；

\item 嵌入成：

\end{enumerate}

\[
Q_k=\operatorname{diag}(I_k,P_k).
\]

\begin{enumerate}

\item 做相似变换：

\end{enumerate}

\[
A\leftarrow Q_k^TAQ_k.
\]

为什么必须左右都乘？因为这是特征值问题，要保持相似，从而保持特征值不变。只左乘会改变特征值，只适合 QR 分解或最小二乘。

\phantomsection

\subsection*{5. QR 迭代}

\addcontentsline{toc}{subsection}{5. QR 迭代}

QR 迭代用于求特征值。第 \code{\detokenize{k}} 步：

\[
A_k=Q_kR_k.
\]

然后：

\[
A_{k+1}=R_kQ_k.
\]

因为：

\[
R_kQ_k=Q_k^TA_kQ_k,
\]

所以 $A_{k+1}$ 与 $A_k$ 正交相似，特征值不变。

在适当条件下，$A_k$ 会趋向上三角或准上三角矩阵，对角线给出特征值。

\phantomsection

\subsection*{6. 位移 QR}

\addcontentsline{toc}{subsection}{6. 位移 QR}

为了加速收敛，可以使用位移：

\[
A_k-\mu_k I=Q_kR_k.
\]

然后：

\[
A_{k+1}=R_kQ_k+\mu_k I.
\]

这仍然与 $A_k$ 相似，因为：

\[
A_{k+1}=Q_k^TA_kQ_k.
\]

位移选择越接近某个特征值，收敛越快。

\phantomsection

\subsection*{7. Hessenberg 结构为什么保持}

\addcontentsline{toc}{subsection}{7. Hessenberg 结构为什么保持}

如果 $A_k$ 是上 Hessenberg，对它做 QR 分解时，$Q_k$ 也是接近 Hessenberg 结构的正交矩阵，乘回：

\[
A_{k+1}=R_kQ_k
\]

仍保持上 Hessenberg。

考试证明不用写太细的实现，只要说明：

\begin{itemize}

\item QR 分解可用局部 Givens 或 Householder 完成；

\item 每步只消次对角以下元素；

\item 乘回后不会产生更低位置的非零元；

\item 所以结构保持。

\end{itemize}

\phantomsection

\subsection*{8. Schur 分解}

\addcontentsline{toc}{subsection}{8. Schur 分解}

复 Schur 分解：

\[
A=QTQ^*,
\]

其中 \code{\detokenize{Q}} 酉，\code{\detokenize{T}} 上三角。\code{\detokenize{T}} 的对角线元素就是 \code{\detokenize{A}} 的特征值。

实 Schur 分解：

\[
A=QTQ^T,
\]

其中 \code{\detokenize{Q}} 正交，\code{\detokenize{T}} 是拟上三角矩阵。拟上三角允许 \code{\detokenize{1}} 阶和 \code{\detokenize{2}} 阶对角块，\code{\detokenize{2}} 阶块对应一对共轭复特征值。

Schur 分解与 QR 迭代的关系：

\begin{itemize}

\item QR 迭代通过正交相似变换逐步逼近 Schur 形；

\item 对称矩阵的 Schur 形可以进一步成为对角阵；

\item 因此对称 QR 最终逼近特征值对角阵。

\end{itemize}

\phantomsection

\subsection*{9. 对称 QR 方法}

\addcontentsline{toc}{subsection}{9. 对称 QR 方法}

若 $A=A^T$，QR 迭代保持对称：

\[
A_{k+1}=Q_k^TA_kQ_k.
\]

相似变换保持对称性：

\[
A_{k+1}^T=A_{k+1}.
\]

对称矩阵可正交对角化，QR 方法通常收敛到对角阵。对称情形还常先化为三对角矩阵，以降低计算量。

\phantomsection

\section*{十、SVD、Moore-Penrose 逆、Jacobi 与二分法}

\addcontentsline{toc}{section}{十、SVD、Moore-Penrose 逆、Jacobi 与二分法}

\phantomsection

\subsection*{1. SVD 定理}

\addcontentsline{toc}{subsection}{1. SVD 定理}

任意矩阵 \code{\detokenize{A}} 都有奇异值分解：

\[
A=U\Sigma V^T.
\]

其中 \code{\detokenize{U,V}} 正交，$\\Sigma$ 是对角型矩阵，非负对角元为奇异值：

\[
\sigma_1\ge \sigma_2\ge \cdots\ge \sigma_r>0.
\]

奇异值和特征值的关系：

\[
\sigma_i=\sqrt{\lambda_i(A^TA)}.
\]

如果 \code{\detokenize{A}} 是方阵且非奇异，则：

\[
\|A\|_2=\sigma_1,
\quad
\|A^{-1}\|_2=\frac{1}{\sigma_n}.
\]

所以：

\[
\kappa_2(A)=\frac{\sigma_1}{\sigma_n}.
\]

\phantomsection

\subsection*{2. 奇异值的变分刻画}

\addcontentsline{toc}{subsection}{2. 奇异值的变分刻画}

常用结论：

\[
\|A\|_2=\max_{\|x\|_2=1}\|Ax\|_2=\sigma_1.
\]

若 \code{\detokenize{A}} 满列秩：

\[
\sigma_1\|x\|_2\ge \|Ax\|_2\ge \sigma_n\|x\|_2.
\]

这个不等式常用于证明最小二乘和条件数相关结论。

\phantomsection

\subsection*{3. Moore-Penrose 广义逆}

\addcontentsline{toc}{subsection}{3. Moore-Penrose 广义逆}

如果：

\[
A=U\Sigma V^T,
\]

则：

\[
A^+=V\Sigma^+U^T.
\]

其中 $\\Sigma^+$ 把非零奇异值取倒数后转置位置。

Moore-Penrose 逆满足四个条件：

\[
AA^+A=A,
\]

\[
A^+AA^+=A^+,
\]

\[
(AA^+)^T=AA^+,
\]

\[
(A^+A)^T=A^+A.
\]

考试常见目标：证明 $AA^+$ 是到 \code{\detokenize{R(A)}} 的正交投影。

只需证明：

\[
(AA^+)^T=AA^+,
\quad
(AA^+)^2=AA^+.
\]

对称且幂等，即正交投影。

\phantomsection

\subsection*{4. Jacobi 特征值方法}

\addcontentsline{toc}{subsection}{4. Jacobi 特征值方法}

Jacobi 方法用于实对称矩阵特征值问题。思想是用一系列平面旋转消去非对角元素。

对称矩阵：

\[
A=A^T.
\]

选择一个非对角元素 $a_{pq}$，构造 Givens 型旋转 \code{\detokenize{J}}，使：

\[
(J^TAJ)_{pq}=0.
\]

不断迭代，非对角元素趋于 0，矩阵趋于对角阵。对角线就是特征值。

它的优点是稳定，适合对称矩阵；缺点是一般不如 QR 方法快。

\phantomsection

\subsection*{5. 二分法与 Sturm 序列}

\addcontentsline{toc}{subsection}{5. 二分法与 Sturm 序列}

对称三对角矩阵的特征值可以用二分法求。核心是：通过 Sturm 序列计算小于某个数 $\mu$ 的特征值个数。

设三对角矩阵：

\[
T=
\begin{pmatrix}
a_1&b_1&&\\
b_1&a_2&b_2&\\
&\ddots&\ddots&\ddots\\
&&b_{n-1}&a_n
\end{pmatrix}.
\]

定义序列：

\[
p_0(\mu)=1,
\]

\[
p_1(\mu)=a_1-\mu,
\]

\[
p_k(\mu)=(a_k-\mu)p_{k-1}(\mu)-b_{k-1}^2p_{k-2}(\mu).
\]

符号变化次数可以用来判断小于 $\mu$ 的特征值个数。然后像普通二分法一样缩小区间。

考试如果只考概念，重点写：

\begin{itemize}

\item 对称三对角矩阵；

\item Sturm 序列计数；

\item 根据计数二分定位特征值。

\end{itemize}

\phantomsection

\section*{十一、证明题统一训练}

\addcontentsline{toc}{section}{十一、证明题统一训练}

\phantomsection

\subsection*{1. 范数证明}

\addcontentsline{toc}{subsection}{1. 范数证明}

所有范数证明都按四步：

\begin{enumerate}

\item 非负性；

\item 零向量判别；

\item 齐次性；

\item 三角不等式。

\end{enumerate}

如果是矩阵范数，再加：

\[
\|AB\|\le \|A\|\|B\|.
\]

对于由已有范数构造的新范数，要尽量把它化回旧范数。例如：

\[
\|A\|_X=\|X^{-1}AX\|_2.
\]

就把所有性质都转移到 $\|.\|_2$。

\phantomsection

\subsection*{2. 收敛证明}

\addcontentsline{toc}{subsection}{2. 收敛证明}

迭代法收敛证明三条路线：

路线一：直接求谱半径。

\[
\rho(B)<1.
\]

路线二：找某个范数。

\[
\|B\|<1\quad \Rightarrow\quad \rho(B)\le \|B\|<1.
\]

路线三：用充分条件。

\begin{itemize}

\item 严格对角占优；

\item SPD；

\item SOR 的 $0<\omega<2$。

\end{itemize}

参数题优先路线一，因为最不容易漏条件。

\phantomsection

\subsection*{3. 正定证明}

\addcontentsline{toc}{subsection}{3. 正定证明}

正定证明三条路线：

\begin{enumerate}

\item 顺序主子式；

\item 特征值；

\item 二次型配方。

\end{enumerate}

如果矩阵带参数，顺序主子式通常最快。

例如：

\[
A=\begin{pmatrix}1&0&a\\0&1&0\\a&0&1\end{pmatrix}
\]

直接看：

\[
\det A=1-a^2.
\]

所以：

\[
|a|<1.
\]

\phantomsection

\subsection*{4. 最小二乘证明}

\addcontentsline{toc}{subsection}{4. 最小二乘证明}

最小二乘证明常从三句话开始：

\begin{enumerate}

\item 最优残差垂直于列空间；

\item 因此 $A^T(b-Ax)=0$；

\item 所以 $A^TAx=A^Tb$。

\end{enumerate}

投影矩阵证明常写：

\[
P^T=P,\quad P^2=P.
\]

只要能证明这两条，就说明 \code{\detokenize{P}} 是正交投影。

\phantomsection

\subsection*{5. 正交变换证明}

\addcontentsline{toc}{subsection}{5. 正交变换证明}

凡是 Householder、Givens、QR、Hessenberg、QR 迭代，都要抓住：

\[
Q^TQ=I.
\]

常见用途：

保持长度：

\[
\|Qx\|_2=\|x\|_2.
\]

保持最小二乘目标：

\[
\|Ax-b\|_2=\|Q^T(Ax-b)\|_2.
\]

保持特征值：

\[
A_{k+1}=Q^TA_kQ.
\]

\phantomsection

\subsection*{6. 共轭梯度证明}

\addcontentsline{toc}{subsection}{6. 共轭梯度证明}

看到 A-共轭，立刻想到：

\[
p_i^TAp_j=0.
\]

线性无关证明：设线性组合为 0，左乘 $p_j^TA$，只剩一项。

逆矩阵展开证明：把任意向量在 A-共轭基下展开，再令 $x=A^{-1}y$。

\phantomsection

\section*{十二、考场作答模板与复习路线}

\addcontentsline{toc}{section}{十二、考场作答模板与复习路线}

\phantomsection

\subsection*{1. 直接法题模板}

\addcontentsline{toc}{subsection}{1. 直接法题模板}

如果题目给矩阵求 LU/PLU：

\begin{enumerate}

\item 写出每步主元；

\item 写出乘子；

\item 写出消元后的 \code{\detokenize{U}}；

\item 把乘子放进 \code{\detokenize{L}}；

\item 若换行，写 \code{\detokenize{P}}；

\item 解方程时一定写 \code{\detokenize{Pb}}。

\end{enumerate}

\phantomsection

\subsection*{2. 最小二乘题模板}

\addcontentsline{toc}{subsection}{2. 最小二乘题模板}

\begin{enumerate}

\item 写目标：

\end{enumerate}

\[
\min_x\|Ax-b\|_2.
\]

\begin{enumerate}

\item 写正交条件：

\end{enumerate}

\[
A^T(Ax-b)=0.
\]

\begin{enumerate}

\item 写正则方程：

\end{enumerate}

\[
A^TAx=A^Tb.
\]

\begin{enumerate}

\item 若要求平方根法，对 $A^TA$ 做 Cholesky。

\item 若要求 QR，写：

\end{enumerate}

\[
Rx=Q_1^Tb.
\]

\phantomsection

\subsection*{3. 迭代法题模板}

\addcontentsline{toc}{subsection}{3. 迭代法题模板}

\begin{enumerate}

\item 写矩阵分裂：

\end{enumerate}

\[
A=D-L-U.
\]

\begin{enumerate}

\item 写迭代矩阵 \code{\detokenize{B}}；

\item 求特征值或谱半径；

\item 判断：

\end{enumerate}

\[
\rho(B)<1.
\]

\begin{enumerate}

\item 若带参数，解不等式。

\end{enumerate}

\phantomsection

\subsection*{4. 正交变换题模板}

\addcontentsline{toc}{subsection}{4. 正交变换题模板}

Householder：

\begin{enumerate}

\item 找要处理的子向量 \code{\detokenize{x}}；

\item 取 $\alpha=\pm \|x\|_2$；

\item 写：

\end{enumerate}

\[
w=\frac{x-\alpha e_1}{\|x-\alpha e_1\|_2}.
\]

\begin{enumerate}

\item 写：

\end{enumerate}

\[
H=I-2ww^T.
\]

Givens：

\begin{enumerate}

\item 写旋转矩阵；

\item 列方程；

\item 加上 $c^2+s^2=1$；

\item 解出 \code{\detokenize{c,s}}。

\end{enumerate}

\phantomsection

\subsection*{5. 特征值综合题模板}

\addcontentsline{toc}{subsection}{5. 特征值综合题模板}

若题目问 Hessenberg：

\begin{itemize}

\item 写 $Q^TAQ=H$；

\item 每步用 Householder 只处理活动子块；

\item 这是正交相似变换，所以特征值不变。

\end{itemize}

若题目问 QR：

\begin{itemize}

\item 写 $A_k=Q_kR_k$；

\item 写 $A_{k+1}=R_kQ_k=Q_k^TA_kQ_k$；

\item 说明相似变换保持特征值。

\end{itemize}

若题目问 Schur：

\begin{itemize}

\item 复 Schur：$A=QTQ^*$，\code{\detokenize{T}} 上三角；

\item 实 Schur：$A=QTQ^T$，\code{\detokenize{T}} 准上三角；

\item QR 迭代逼近 Schur 形。

\end{itemize}

\phantomsection

\subsection*{6. 最可能考题清单}

\addcontentsline{toc}{subsection}{6. 最可能考题清单}

最高优先级：

\begin{enumerate}

\item 证明相似诱导范数是矩阵范数；

\item 给 3 阶矩阵做 LU/PLU 并解方程；

\item 最小二乘正则方程加 Cholesky；

\item Householder 与 Givens 手算；

\item 参数矩阵正定性与 Jacobi/G-S 收敛；

\item Jacobi 迭代矩阵和谱半径；

\item A-共轭向量线性无关与 $A^{-1}$ 展开；

\item Hessenberg、QR、Schur、SVD 理论综合。

\end{enumerate}

中等优先级：

\begin{enumerate}

\item Frobenius 范数和正交不变性；

\item 投影矩阵 $P=A(A^TA)^{-1}A^T$；

\item 追赶法；

\item 条件数与残量误差；

\item 幂法、逆迭代、位移 QR；

\item Jacobi 特征值方法与二分法。

\end{enumerate}

\phantomsection

\subsection*{7. 最后一天怎么复习}

\addcontentsline{toc}{subsection}{7. 最后一天怎么复习}

最后一天不要重新看所有细节，按题型刷：

\begin{enumerate}

\item 手算一遍 LU/PLU；

\item 手算一遍 Cholesky 或 LDLT；

\item 手算一遍最小二乘正则方程；

\item 手算一遍 Householder/Givens；

\item 写出 Jacobi/G-S/SOR 三个迭代矩阵；

\item 默写 CG 递推；

\item 默写 QR 迭代与 Schur 分解；

\item 背 5 个证明模板：范数、正定、收敛、投影、A-共轭。

\end{enumerate}

\phantomsection

\subsection*{8. 最容易救分的卷面习惯}

\addcontentsline{toc}{subsection}{8. 最容易救分的卷面习惯}

\begin{itemize}

\item 先声明记号约定，例如 $PA=LU$；

\item 每个算法题写清中间变量，不只写最终答案；

\item 证明题先写目标性质，再逐条验证；

\item 有参数就把“当且仅当”的范围写完整；

\item 最小二乘题一定写残差垂直；

\item 迭代题一定写谱半径；

\item 特征值题一定写相似变换；

\item 正交变换题一定写 $Q^TQ=I$。

\end{itemize}

\end{document}
